% tbsrc/summary.tex — Summary and the Example Ladder
% Status: written (first draft). Recap of the Part I arc around the data-model
% spine; the example-ladder table (physics rows) and the distinct solver axis;
% forward pointer to the remaining backgrounds.
\section{Summary and the Example Ladder}
\label{sec:summary}

Part~I built one equation and the algorithm that inverts it. The mean counts on
LOR $i$,
\begin{equation}
  \pred_i \;=\; n_i\,a_i\,\bigl(\Aop\,\Gop\,\xv\bigr)_i \;+\; s_i + r_i ,
  \label{eq:sum-spine}
\end{equation}
is the spine that every section hangs from: the geometry and grid fix what $\Aop$
means (Sections~\ref{sec:geometry}--\ref{sec:projection}), the resolution is
$\Gop$ (Section~\ref{sec:resolution}), the physical effects are $n,a,s,r$
(Section~\ref{sec:datamodel}), the Poisson statistics turn \eqref{eq:sum-spine}
into the likelihood (Section~\ref{sec:statistics}), and MLEM/OSEM minimize the
resulting objective (Sections~\ref{sec:mlem}--\ref{sec:acceleration}), assembled
into the two-phase pipeline of Section~\ref{sec:pipeline}. The library mirrors the
split: the multiplicative factors enter as $\code{mult}=\nrm\odot\att$, the
additive backgrounds as $\code{contamination}=\sct+\rnd$, the resolution is
injected at simulation time and left in the data, and the reconstruction needs
only these plus the sensitivity image $\sensv=\Aadj(\code{mult})$.

\subsection{The example ladder}
\label{sec:summary:ladder}

The worked examples (Part~II) switch the terms of \eqref{eq:sum-spine} on one at a
time, each isolating a single effect against everything before it already
working. Table~\ref{tab:ladder} is that ladder.

\begin{table}[h]
  \centering
  \begin{tabular}{@{}llllll@{}}
    \toprule
    Example & $\Gop$ & $\att$ & $\rnd$ & $\sct$ & what it isolates \\
    \midrule
    Attenuation & $\mathbf I$ & water $\mu$-map & $0$ & $0$ & cupping and its correction \\
    Resolution  & $5\,\mathrm{mm}$ & water & $0$ & $0$ & the blur $\Gop$, recovery of $\Gop\xv$ \\
    $+$ Randoms & $5\,\mathrm{mm}$ & water & flat & $0$ & a smooth additive background \\
    $+$ Scatter & $5\,\mathrm{mm}$ & water & flat & SSS/MC & object-dependent background \\
    \bottomrule
  \end{tabular}
  \caption{The physics example ladder. Each row adds one term of
    \eqref{eq:sum-spine} to the row above; normalization is $\nrm=\bm1$
    throughout (the continuous detector's geometric sensitivity is already in
    $\sensv$, Section~\ref{sec:datamodel}). The first two examples exist; randoms
    and scatter follow.}
  \label{tab:ladder}
\end{table}

The ladder is cumulative by design: because the reconstruction code is fixed
(Section~\ref{sec:pipeline}), turning on a term changes only how the data are
simulated, so each example's new effect is the \emph{only} thing that differs from
its predecessor and its consequence is unambiguous.

\subsection{A second axis: the solver}
\label{sec:summary:solver}

One worked example does not sit on this ladder. The acceleration example holds the
physics fixed --- it reuses the resolution phantom and attenuation --- and varies
the \emph{solver} instead, comparing MLEM with OSEM and exhibiting
semi-convergence. This is a genuinely different axis: the ladder switches on terms
of the forward model \eqref{eq:sum-spine}, while the solver study switches on
properties of the inversion (subset acceleration, the limit cycle, early
stopping). The two are orthogonal, which is exactly why the same generic
reconstruction runs every row of the ladder and every solver in the study.

\subsection{What follows}
\label{sec:summary:next}

The ladder's open rows are the additive backgrounds. Randoms add a smooth,
activity-independent term estimated from a delayed window; scatter adds a broad,
object-dependent term from an external single-scatter or Monte-Carlo model
(Section~\ref{sec:datamodel}). Both enter through \code{contamination} without
touching the reconstruction, completing the data model \eqref{eq:sum-spine} on
real-scale data. With geometry, resolution, attenuation, statistics, and the
solver in hand, those examples are the last step from the idealized forward model
to a quantitative reconstruction~\cite{crysp2026}.
